\MathBookSourcePage{235}
\subsection{引理~\ref{lem:ch08-weighted-adjoint-regularity} 和引理~\ref{lem:ch08-weighted-primal-regularity} 的证明}
\label{sec:ch08-proofs-weighted-regularity-lemmas}
\setcounter{thmcount}{0}
\setcounter{equation}{0}

先考虑引理~\ref{lem:ch08-weighted-adjoint-regularity}. 由 Hölder 不等式、式~\eqref{eq:ch08-primal-regularity} 和习题~\ref{exr:ch08-12},
\begin{align*}
\MathBookFormulaID{formula-ch08-adjoint-regularity-holder-chain}
\int_\Omega\sigma^{-d-\lambda}|\nabla_2v|^2\,dx
&\leq\left(\int_\Omega\sigma^{-(d+\lambda)p'}\,dx\right)^{1/p'}
\|\nabla_2v\|_{L^{2p}(\Omega)}^2\\
&\leq C\zeta^{-d-\lambda+d/p'}\|v\|_{W_{2p}^2(\Omega)}^2
=C\zeta^{-\lambda-d/p}\|v\|_{W_{2p}^2(\Omega)}^2\\
&\leq C\zeta^{-\lambda-d/p}\|f\|_{L^{2p}(\Omega)}^2
\leq C\zeta^{-\lambda-d/p}
\|\nabla f\|_{L^{2pd/(d+2p)}(\Omega)}^2.
\end{align*}
若 $p<d/(d-2)$, 取 $q=(d+2p)/(pd)$, 再次使用 Hölder 不等式可得
\[
\MathBookFormulaID{formula-ch08-adjoint-regularity-weighted-holder}
\|\nabla f\|_{L^{2pd/(d+2p)}(\Omega)}^2
\leq\left(\int_\Omega\sigma^{-(4-d-\lambda)q/(q-1)}\,dx\right)^{(q-1)/q}
\left(\int_\Omega\sigma^{4-d-\lambda}|\nabla f|^2\,dx\right)^{1/q}.
\]
由 $\sigma$ 的定义,
\[
\MathBookFormulaID{formula-ch08-adjoint-weight-integral}
\left(\int_\Omega\sigma^{-(4-d-\lambda)q/(q-1)}\,dx\right)^{q-1}
\leq C\zeta^{-(4-d-\lambda)+d(q-1)}
=C\zeta^{-2+\lambda+d/p}.
\]
因为假设给出 $d/(2-\lambda)<\mu$, 可选择满足这些约束的 $p$. 合并上述估计即证明引理~\ref{lem:ch08-weighted-adjoint-regularity}.

现在考虑引理~\ref{lem:ch08-weighted-primal-regularity}. 展开 $\nabla_2(\sigma^{(d+\lambda)/2}v)$, 得
\MBSourceItem{1}
\begin{equation}
\MathBookFormulaID{formula-ch08-weighted-second-derivative-expansion}
\label{eq:ch08-weighted-second-derivative-expansion}
\sigma^{d+\lambda}|\nabla_2v|^2
\leq\left|\nabla_2\bigl(\sigma^{(d+\lambda)/2}v\bigr)\right|^2
+C\bigl(\sigma^{d+\lambda-2}|\nabla v|^2+\sigma^{d+\lambda-4}v^2\bigr).
\end{equation}
为使用式~\eqref{eq:ch08-primal-regularity}, 考察 $\sigma^{(d+\lambda)/2}v$ 所满足的方程. 对 $\phi\in V$, 分部积分并利用系数导数展开可写成
\[
\MathBookFormulaID{formula-ch08-weighted-coefficient-equation}
a\bigl(\phi,\sigma^{(d+\lambda)/2}v\bigr)=(F,\phi),
\]
其中
\[
\MathBookFormulaID{formula-ch08-weighted-coefficient-forcing}
F=\sigma^{(d+\lambda)/2}\nu\cdot\nabla f
-\sum_{i,j=1}^{d}\bigl(a_{ij}(\sigma^{(d+\lambda)/2})_{,j}v\bigr)_{,i}
-\sum_{i,j=1}^{d}a_{ij}(\sigma^{(d+\lambda)/2})_{,j}v_{,i}
-\sum_i b_i(\sigma^{(d+\lambda)/2})_{,i}v.
\]
因此, 由式~\eqref{eq:ch08-weight-element-comparability}、式~\eqref{eq:ch08-weight-derivative-bound} 和式~\eqref{eq:ch08-primal-regularity},
\[
\MathBookFormulaID{formula-ch08-weighted-primal-preliminary-bound}
\begin{aligned}
\int_\Omega\sigma^{d+\lambda}|\nabla_2v|^2\,dx
&\leq C\left(\int_\Omega\sigma^{d+\lambda}|\nabla f|^2\,dx
+\int_\Omega\sigma^{d+\lambda-2}|\nabla v|^2\,dx
+\int_\Omega\sigma^{d+\lambda-4}v^2\,dx\right)\\
&\quad+C\int_\Omega A(x)^2\sigma^{d+\lambda-2}v^2\,dx.
\end{aligned}
\]
其中 $A(x)=\max_{ij}|a_{ij,i}(x)|$. 对最后一个积分应用 Hölder 和 Sobolev 不等式. 当 $d=2$ 时要求 $p>2$, 当 $d=3$ 时要求 $p\geq12/5$, 即得到
\[
\MathBookFormulaID{formula-ch08-weighted-primal-reduced-bound}
\int_\Omega\sigma^{d+\lambda}|\nabla_2v|^2\,dx
\leq C\left(\int_\Omega\sigma^{d+\lambda}|\nabla f|^2\,dx
+\int_\Omega\sigma^{d+\lambda-2}|\nabla v|^2\,dx
+\int_\Omega\sigma^{d+\lambda-4}v^2\,dx\right).
\]
回顾式~\eqref{eq:ch08-dual-weighted-error-problem} 并约定 $v_{,0}=v$, 有
\MBSourceItem{2}
\begin{align}
\MathBookFormulaID{formula-ch08-weighted-primal-energy-identity}
C_a\int_\Omega\sigma^{d+\lambda-2}|\nabla v|^2\,dx
&\leq\int_\Omega\sigma^{d+\lambda-2}
\sum_{i,j=1}^{d}a_{ij}v_{,i}v_{,j}\,dx\notag\\
&=a(v,\sigma^{d+\lambda-2}v)
-\int_\Omega\sum_{i,j=1}^{d}a_{ij}v_{,i}v
(\sigma^{d+\lambda-2})_{,j}\,dx
+\int_\Omega\sum_{i=0}^{d}b_iv_{,i}v\sigma^{d+\lambda-2}\,dx\notag\\
&=((\nu\cdot\nabla f),\sigma^{d+\lambda-2}v)
-\int_\Omega\sum_{i,j=1}^{d}a_{ij}v_{,i}v
(\sigma^{d+\lambda-2})_{,j}\,dx\notag\\
&\quad+\int_\Omega\sum_{i=0}^{d}b_iv_{,i}v\sigma^{d+\lambda-2}\,dx.
\label{eq:ch08-weighted-primal-energy-identity}
\end{align}
与前面相同地估计, 得
\MBSourceItem{3}
\begin{equation}
\MathBookFormulaID{formula-ch08-weighted-primal-energy-bound}
\label{eq:ch08-weighted-primal-energy-bound}
\int_\Omega\sigma^{d+\lambda-2}|\nabla v|^2\,dx
\leq C\left(\int_\Omega\sigma^{d+\lambda}|\nabla f|^2\,dx
+\int_\Omega\sigma^{d+\lambda-4}v^2\,dx\right).
\end{equation}
Hölder 不等式给出
\[
\MathBookFormulaID{formula-ch08-weighted-primal-l2-holder}
\int_\Omega\sigma^{d+\lambda-4}v^2\,dx
\leq C\zeta^{(d+\lambda-4)+d/P'}\|v\|_{L^{2P}(\Omega)}^2,
\]
其中 $P'>d/(4-\lambda-d)$. 再作一次对偶论证: 令 $w$ 解式~\eqref{eq:ch08-primal-problem}, 右端为 $\operatorname{sign}(v)|v|^{2P-1}$. 则
\begin{align*}
\MathBookFormulaID{formula-ch08-duality-l2p-chain}
\|v\|_{L^{2P}(\Omega)}^{2P}
&=(\operatorname{sign}(v)|v|^{2P-1},v)=a(w,v)\\
&=(\nu\cdot\nabla f,w)=-(f,\nu\cdot\nabla w)\\
&\leq\|f\|_{L^r(\Omega)}\|w\|_{W_{r'}^1(\Omega)}
\leq C\|f\|_{L^r(\Omega)}\|w\|_{W_{2P/(2P-1)}^2(\Omega)}\\
&\leq C\|f\|_{L^r(\Omega)}
\bigl\||v|^{2P-1}\bigr\|_{L^{2P/(2P-1)}(\Omega)}
=C\|f\|_{L^r(\Omega)}\|v\|_{L^{2P}(\Omega)}^{2P-1}.
\end{align*}
因此
\[
\MathBookFormulaID{formula-ch08-duality-l2p-bound}
\|v\|_{L^{2P}(\Omega)}\leq C\|f\|_{L^r(\Omega)},
\qquad r=\frac{2Pd}{2P+d}.
\]
Sobolev 不等式要求
\[
\MathBookFormulaID{formula-ch08-duality-exponent-relation}
\frac1{r'}=\frac{2P-1}{2P}-\frac1d
=1-\frac1{2P}-\frac1d.
\]
条件 $P'>d/(4-\lambda-d)$ 等价于
\[
\MathBookFormulaID{formula-ch08-duality-upper-r-condition}
r<\frac{2d}{2d-2+\lambda}.
\]
为应用式~\eqref{eq:ch08-primal-regularity}, 还必须有
\[
\MathBookFormulaID{formula-ch08-duality-lower-r-condition}
r>\left(1-\frac1\mu+\frac1d\right)^{-1}.
\]
若 $\mu>d$, 后一条件成为空条件; 若 $\mu>2d/(4-\lambda)$, 则存在满足两个条件的 $r$ 的开区间. 因而 Hölder 不等式给出
\begin{align*}
\MathBookFormulaID{formula-ch08-weighted-primal-lower-order-bound}
\int_\Omega\sigma^{d+\lambda-4}v^2\,dx
&\leq C\zeta^{2d(1-1/r)+\lambda-2}\|f\|_{L^r(\Omega)}^2\\
&\leq C\zeta^{2d(1-1/r)+\lambda-2}
\int_\Omega\sigma^{d+\lambda}f^2\,dx
\left(\int_\Omega\sigma^{-(d+\lambda)r/(2-r)}\,dx\right)^{(2-r)/r}\\
&\leq C\zeta^{-2}\int_\Omega\sigma^{d+\lambda}f^2\,dx.
\end{align*}
代回式~\eqref{eq:ch08-weighted-primal-energy-bound} 及前面的二阶导数估计, 完成引理~\ref{lem:ch08-weighted-primal-regularity} 的证明.

\MBSourceItem{4}
\begin{Remark}
\label{rem:ch08-index-engineering}
上述证明反复使用带有初看似乎很神奇的指标的 Hölder 不等式. 这种论证技巧可称为“指标工程”. 其基本科学原则是量纲分析: 引理~\ref{lem:ch08-weighted-adjoint-regularity} 中不等式的两端具有相同的量纲. 若 $h$ 的长度单位为 $L$, 则左端的单位为 $L^{2-d-\lambda}$, 因为 $\zeta$ 的单位为 $L$; 右端同样如此. 这类检查不会给出精确估计, 但能判断推导是否可能成立.
\end{Remark}
